\documentclass[11pt]{article}
\usepackage[T1]{fontenc}
\usepackage[margin=1in]{geometry}
\usepackage{amsmath,amssymb,amsthm,microtype,graphicx}
\usepackage[colorlinks=true,urlcolor=blue,citecolor=blue]{hyperref}
\newtheorem{theorem}{Theorem}
\newtheorem{corollary}[theorem]{Corollary}
\newcommand{\E}{\mathbb E}
\newcommand{\Rp}{R_p}
\newcommand{\Mp}{M_p}
\newcommand{\op}{\mathrm{op}}
\title{Sharp sparse Rademacher comparison for orthogonal monochromatic blocks}
\author{Candidate partial result for \texttt{fa\_banach\_001}}
\date{11 August 2026}

\begin{document}
\maketitle

\noindent\textbf{Status.}
Candidate substantial partial result; likely valid.  This is not a solution
of the full arbitrary-matrix or circulant problem.

\section{Source question}

For a scalar matrix $A=(a_{ij})$ and $p\geq1$, Lata\l a and
\'Swi\k atkowski define
\[
 \Rp(A):=\sup_{\|s\|_2,\|t\|_2\leq1}
 \left\|\sum_{i,j}a_{ij}\varepsilon_{ij}s_it_j\right\|_p
\]
and prove in Proposition 1.5 that
\[
 \frac12\Mp(A)\leq\Rp(A)
 \lesssim\sqrt{\log(p+1)}\,\Mp(A),
 \qquad
 \Mp(A):=\max_{|I|\leq p}
 \bigl\|(|a_{ij}|\mathbf1_{(i,j)\in I})\bigr\|_{\op}.
\]
They then ask whether the logarithmic factor is necessary
\cite[Proposition 1.5 and the following sentence, p.~4]{LS2021}.

\begin{figure}[ht]
 \centering
 \includegraphics[width=0.96\textwidth]{figures/open_problem_crop.png}
 \caption{The source question on arXiv PDF page 4.}
\end{figure}

The same paper proves the constant-factor estimate
\begin{equation}\label{eq:zeroone}
 \Rp(B)\leq C_0\Mp(B)
\end{equation}
for every $0$--$1$ matrix $B$, with a universal constant $C_0$
\cite[Proposition 1.4]{LS2021}.  The result below shows that this sharp class
has a dimension-free direct-sum closure that permits arbitrarily many
different coefficient scales.

\section{Proof intuition}

An orthogonal block decomposition splits both test vectors into square-summable
pieces.  The source theorem controls the Rademacher sum inside each block by
the global sparse norm times the product of the two corresponding block
lengths.  Minkowski adds these products, while Cauchy--Schwarz bounds their
sum by one.  Thus neither the number of blocks nor their coefficient scales
can create the logarithmic accumulation seen in a generic dyadic
decomposition.  The unresolved case is exactly where different scales share
rows or columns, destroying this two-sided square summability.

\section{Direct-sum stability}

We allow rectangular matrices in the statement; they can equivalently be
padded by zero rows or columns.

\begin{theorem}[orthogonal direct-sum stability]\label{thm:stability}
Fix $p\geq1$.  Suppose $A$ is, after independent permutations of its rows and
columns, an orthogonal direct sum
\[
 A=A_1\oplus\cdots\oplus A_m.
\]
If a constant $C$ satisfies
\[
 \Rp(A_\ell)\leq C\Mp(A_\ell)
 \qquad(1\leq\ell\leq m),
\]
then the same constant satisfies
\[
 \Rp(A)\leq C\Mp(A).
\]
\end{theorem}

\begin{proof}
Fix $s,t$ with $\|s\|_2,\|t\|_2\leq1$.  Let $s_\ell,t_\ell$ be their
restrictions to the row and column spaces of $A_\ell$, and put
\[
 \alpha_\ell=\|s_\ell\|_2,\qquad
 \beta_\ell=\|t_\ell\|_2.
\]
When these quantities are nonzero, write
$s_\ell=\alpha_\ell u_\ell$ and
$t_\ell=\beta_\ell v_\ell$ with $u_\ell,v_\ell$ unit vectors.  A zero
restriction simply contributes zero.  Minkowski's inequality and the assumed
block estimates give
\begin{align*}
 \left\|\sum_{i,j}a_{ij}\varepsilon_{ij}s_it_j\right\|_p
 &\leq \sum_{\ell=1}^m
 \alpha_\ell\beta_\ell
 \left\|\sum_{(i,j)\text{ in block }\ell}
 (A_\ell)_{ij}\varepsilon_{ij}(u_\ell)_i(v_\ell)_j\right\|_p\\
 &\leq C\sum_{\ell=1}^m
 \alpha_\ell\beta_\ell\Mp(A_\ell).
\end{align*}
Every set of at most $p$ entries in $A_\ell$ is also an admissible set in
$A$, and embedding the resulting sparse matrix does not change its operator
norm.  Hence $\Mp(A_\ell)\leq\Mp(A)$.  Since the row blocks and column blocks
are respectively disjoint,
\[
 \sum_\ell\alpha_\ell^2\leq\|s\|_2^2\leq1,
 \qquad
 \sum_\ell\beta_\ell^2\leq\|t\|_2^2\leq1.
\]
Cauchy--Schwarz now yields
\[
 \sum_\ell\alpha_\ell\beta_\ell\leq1,
\]
so the displayed Rademacher norm is at most $C\Mp(A)$.  Taking the supremum
over $s,t$ proves the assertion.
\end{proof}

\begin{corollary}[orthogonal monochromatic blocks]\label{cor:blocks}
Let $R_\ell$ be pairwise disjoint row sets and $C_\ell$ pairwise disjoint
column sets.  For each $\ell$, let
$E_\ell\subset R_\ell\times C_\ell$ and let $\lambda_\ell$ be a scalar.
Suppose that
\[
 a_{ij}=\lambda_\ell\mathbf1_{E_\ell}(i,j)
 \quad\text{on }R_\ell\times C_\ell,
 \qquad a_{ij}=0\quad\text{otherwise}.
\]
Then, with universal constants independent of all the data,
\[
 \Rp(A)\asymp\Mp(A).
\]
In particular, the factor $\sqrt{\log(p+1)}$ is unnecessary on this class,
even when the number of blocks and the dynamic range of
$|\lambda_\ell|$ are arbitrary.
\end{corollary}

\begin{proof}
The law defining $\Rp$ is unchanged when deterministic entry signs are
absorbed into the independent Rademacher variables.  Thus each nonzero block
may be written as $|\lambda_\ell|B_\ell$ with $B_\ell$ a $0$--$1$ matrix.
Homogeneity and \eqref{eq:zeroone} give
\[
 \Rp(|\lambda_\ell|B_\ell)
 \leq C_0\Mp(|\lambda_\ell|B_\ell).
\]
Theorem~\ref{thm:stability} supplies the upper estimate for $A$.

For completeness, the reverse estimate is the general lower bound from the
source.  Indeed, choose $I$, $|I|\leq p$, and nonnegative unit vectors $u,v$
that norm the nonnegative sparse matrix
$(|a_{ij}|\mathbf1_{(i,j)\in I})$.  Conditional expectation onto the signs in
$I$ shows that deleting all other summands cannot increase the $L_p$ norm.
On the event that the signs in $I$ align with the deterministic coefficient
signs, which has probability $2^{-|I|}$, the selected sum equals
$\Mp(A)$.  Therefore its $L_p$ norm is at least
$2^{-|I|/p}\Mp(A)\geq\frac12\Mp(A)$.  Taking the optimal $I$ completes the
proof.
\end{proof}

\section{Verification, scope, and novelty status}

The proof has no computational dependency.  The key normalization identities,
the local-to-global inequality $\Mp(A_\ell)\leq\Mp(A)$, the handling of zero
blocks and signs, and the noninteger-$p$ lower-bound constant were separately
audited in \texttt{verification\_report.md}.

The result does not cover a decomposition whose pieces merely have disjoint
entries: disjointness is required simultaneously in their row spaces and in
their column spaces.  In particular, the permutation diagonals of a circulant
matrix overlap in every row and column, so this theorem does not settle the
main circulant conjecture.  The obstruction persists in current work:
arXiv:2405.13656 improves the general and circulant estimates but explicitly
states that a simple sharp estimate for $R_A(p)$ remains unclear
\cite{Latala2024}; arXiv:2512.23673 still treats the broader conjecture as open
\cite{Meller2025}.

The run indexes and these two later arXiv sources were searched for the exact
sparse comparison, block-diagonal variants, and direct-sum terminology.  No
exact statement of Theorem~\ref{thm:stability} or
Corollary~\ref{cor:blocks} was found.  Nonetheless the derivation is elementary
from Proposition 1.4 of the source, so novelty confidence is modest.  The
recommended human review is to verify the scope and decide whether this
closure lemma is useful enough to retain as a substantial partial result.

\begin{thebibliography}{9}
\bibitem{LS2021}
R.~Lata\l a and W.~\'Swi\k atkowski,
\emph{Norms of Randomized Circulant Matrices},
arXiv:2106.03139,
\url{https://arxiv.org/abs/2106.03139}.

\bibitem{Latala2024}
R.~Lata\l a,
\emph{On the spectral norm of Rademacher matrices},
arXiv:2405.13656,
\url{https://arxiv.org/abs/2405.13656}.

\bibitem{Meller2025}
R.~Meller,
\emph{Spectral norm of matrices with independent entries up to polyloglog},
arXiv:2512.23673,
\url{https://arxiv.org/abs/2512.23673}.
\end{thebibliography}

\end{document}
